% ----------------------
\subsection{Section 66 (29 October 1831)}
% ----------------------
	\label{DC66}
	
	% -----
	\subsubsection{Historical Background}
	% -----
		\label{DC66-HB}
		
		% --
		\paragraph{JSP}
		% --
		
			``On 29 October 1831, William E. McLellin wrote in his journal, `The Lord condecended to hear my prayr and give me a revelation of his will, through his prophet or Seer (Joseph).' McLellin, a recent convert from Paris, Illinois, met JS for the first time at the 25--26 October 1831 conference in Orange, Cuyahoga County, Ohio, where McLellin was ordained to the high priesthood. At the conclusion of the conference, he accompanied JS to Hiram, Ohio, arriving there on 29 October. McLellin later recalled that on that day, he `went before the Lord in secret, and on my knees asked him to reveal the answer to five questions through his Prophet.' At McLellin's request, JS dictated a revelation for him, perhaps in the southeast upstairs bedroom of the John and Alice (Elsa) Jacobs Johnson home, where JS worked on his revision of the Bible. According to McLellin, the revelation answered the questions `to my full and entire satisfaction.' Although McLellin never explained what his five queries were, the revelation's contents indicate that he was probably concerned about his standing before God and about what the Lord desired him to do.
			
			``McLellin recounted that he wrote the words of this revelation as JS spoke them. Two copies of the revelation in McLellin's handwriting exist, but it does not appear that either is the original manuscript. One copy is in McLellin's journal, probably made soon after the revelation was dictated. The other is the copy in McLellin's notebook, featured below. McLellin apparently inscribed this copy sometime before 16 November 1831, when he departed on a mission. Three revelations precede the 29 October revelation in McLellin's notebook, including one dated 30 October 1831, indicating McLellin did not make these copies before that date. McLellin's journal corroborates this dating, stating that he stayed in Hiram from 29 October to 16 November and `read and copyed revelations, \&c.' John Whitmer also copied the revelation into Revelation Book 1 sometime before he left for Missouri on 20 November.
			
			``McLellin's notebook copy appears to be a more complete reflection of the original revelation than either the journal copy or the copy John Whitmer made in Revelation Book 1. The journal copy probably predates the other two versions, but the spelling of certain words (`shew' instead of `show,' which is in the journal copy, for example), the use of contractions, and the lack of punctuation suggest that the copies in McLellin's notebook and in Revelation Book 1 were derived from a nonextant copy of the original. The notebook version also appears to be more complete than the copy in Revelation Book 1 because the notebook contains an endnote regarding McLellin belonging to the lineage of Ephraim in the Old Testament (a note that also appears in the journal copy).''
			
		% --
		\paragraph{HDDC} 
		% --
		
			``One of the most interesting sidelights of this revelation occurred after William E. McLellin apostatized. Elder Heber C. Kimball related how a warning given in this revelation was literally fulfilled. The following is taken from Heber C. Kimball's journal.
			
			\begin{quote}
				William E. McLellin wanted to know where Heber C. Kimball was, some one pointed me out to him, as I was sitting on the ground: he came up to me and said `Bro. Heber what do you think of Joseph Smith the fallen prophet, now, has he not led you blindfolded long enough; look and see yourself poor, your family stripped and robbed, and your brethren in the same fix, are you satisfied with Joseph? I replied `Yes, I am more satisfied with him a hundred fold, then ever I was before, for I see you in the very position that he foretold you would be in; a Judas to betray your brethren, if you did not forsake your adultery, fornication, lying and abominations. Where are you? What are you about? you, and Hinkle, and scores of others; have you not betrayed Joseph and his brethren into the hands of the Mob, as Judas did Jesus? Yes, verily, you have; I tell you Mormonism is true, and Joseph is a true prophet of the living God, and you with all others that turn therefrom will be damned and go to hell, and Judas will rule over you.%
					\footnote{%
						% \label{}%
						HDDC note: ``Heber C. Kimball Journal, Book 94C, p. 88, located in the HDC.''
						}
			\end{quote}
			
			``William E. McLellin never did repent in order to affiliate with the Church again and finally died an obscure person in Independence, Missouri on April 24, 1883.''
			
	% -----
	\subsubsection{Versions}
	% -----
		\label{DC66-V}
		
		See the \href{https://drive.google.com/file/d/1goAvNz8jS4R8slYfMG_db55Ty2z_vmgK/view?usp=sharing}{sections comparison} document.
	
		\begin{compactdesc}
			\item[JSP] \hyperref[EMS-1831-JS-29-OCT-1831]{29 October 1831}
			\item[BC] ---
			\item[DC 1835] Section 74, 203--204
			\item[TS] 5:7 (1 April 1844), 482--483
			\item[DC 1844] Section 75, 310--312
			\item[MS] 5:11 (April 1845), 164
		\end{compactdesc}

	% -----
	\subsubsection{Section 66:1} % *DC66-1 
	% -----
		\label{DC66-1}

		BEHOLD, thus saith the Lord unto my servant William E. McLellin---Blessed are you, inasmuch as you have turned away from your iniquities, and have received my truths, saith the Lord your Redeemer, the Savior of the world, even of as many as believe on my name.

		% \begin{framed}
		% \scriptsize
			% \begin{compactdesc}
				% \item[JSP] 
				% % \item[BC] 
				% % \item[]
			% \end{compactdesc}
		% \end{framed}

	% -----
	\subsubsection{Section 66:2} % *DC66-2 
	% -----
		\label{DC66-2}

		Verily I say unto you, blessed are you for receiving mine everlasting covenant, even the fulness of my gospel, sent forth unto the children of men, that they might have life and be made partakers of the glories which are to be revealed in the last days, as it was written by the prophets and apostles in days of old.

		% \begin{framed}
		% \scriptsize
			% \begin{compactdesc}
				% \item[JSP] 
				% % \item[BC] 
				% % \item[]
			% \end{compactdesc}
		% \end{framed}

	% -----
	\subsubsection{Section 66:3} % *DC66-3 
	% -----
		\label{DC66-3}

		Verily I say unto you, my servant William, that you are clean, but not all; repent, therefore, of those things which are not pleasing in my sight, saith the Lord, for the Lord will show them unto you.

		% \begin{framed}
		% \scriptsize
			% \begin{compactdesc}
				% \item[JSP] 
				% % \item[BC] 
				% % \item[]
			% \end{compactdesc}
		% \end{framed}

	% -----
	\subsubsection{Section 66:4} % *DC66-4 
	% -----
		\label{DC66-4}

		And now, verily, I, the Lord, will show unto you what I will concerning you, or what is my will concerning you.

		% \begin{framed}
		% \scriptsize
			% \begin{compactdesc}
				% \item[JSP] 
				% % \item[BC] 
				% % \item[]
			% \end{compactdesc}
		% \end{framed}

	% -----
	\subsubsection{Section 66:5} % *DC66-5 
	% -----
		\label{DC66-5}

		Behold, verily I say unto you, that it is my will that you should proclaim my gospel from land to land, and from city to city, yea, in those regions round about where it has not been proclaimed.

		% \begin{framed}
		% \scriptsize
			% \begin{compactdesc}
				% \item[JSP] 
				% % \item[BC] 
				% % \item[]
			% \end{compactdesc}
		% \end{framed}

	% -----
	\subsubsection{Section 66:6} % *DC66-6 
	% -----
		\label{DC66-6}

		Tarry not many days in this place; go not up unto the land of Zion as yet; but inasmuch as you can send, send; otherwise, think not of thy property.

		% \begin{framed}
		% \scriptsize
			% \begin{compactdesc}
				% \item[JSP] 
				% % \item[BC] 
				% % \item[]
			% \end{compactdesc}
		% \end{framed}

	% -----
	\subsubsection{Section 66:7} % *DC66-7 
	% -----
		\label{DC66-7}

		Go unto the eastern lands, bear testimony in every place, unto every people and in their synagogues, \hyperref[REASON]{reasoning} with the people.

		% \begin{framed}
		% \scriptsize
			% \begin{compactdesc}
				% \item[JSP] 
				% % \item[BC] 
				% % \item[]
			% \end{compactdesc}
		% \end{framed}

	% -----
	\subsubsection{Section 66:8} % *DC66-8 
	% -----
		\label{DC66-8}

		Let my servant Samuel H. Smith go with you, and forsake him not, and give him thine instructions; and he that is faithful shall be made strong in every place; and I, the Lord, will go with you.

		% \begin{framed}
		% \scriptsize
			% \begin{compactdesc}
				% \item[JSP] 
				% % \item[BC] 
				% % \item[]
			% \end{compactdesc}
		% \end{framed}

	% -----
	\subsubsection{Section 66:9} % *DC66-9 
	% -----
		\label{DC66-9}

		Lay your hands upon the sick, and they shall recover.  Return not till I, the Lord, shall send you.  Be patient in affliction.  Ask, and ye shall receive; knock, and it shall be opened unto you.

		% \begin{framed}
		% \scriptsize
			% \begin{compactdesc}
				% \item[JSP] 
				% % \item[BC] 
				% % \item[]
			% \end{compactdesc}
		% \end{framed}

	% -----
	\subsubsection{Section 66:10} % *DC66-10 
	% -----
		\label{DC66-10}

		Seek not to be cumbered.%
			\footnote{%
				% \label{}%
				Forms of the word ``cumber'' appear eight times in the scriptures; this is the only DC use. There are five uses in \hyperref[JACOB-5]{Jacob 5} and then two in Luke, \hyperref[LUKE-10-40]{10:40} and \hyperref[LUKE-13-7]{13:7}. Here are the 1828 definitions:
					\begin{compactitem}
						\item[1] To load; to crowd.
						\item[2] To check, stop or retard, as by a load or weight; to make motion difficult; to obstruct.
						\item[3] To perplex or embarrass; to distract or trouble.
						\item[4] To trouble; to be troublesome to; to cause trouble or obstruction in, as any thing useless.
					\end{compactitem}
				}
		Forsake all unrighteousness.  Commit not adultery---a temptation with which thou hast been troubled.

		% \begin{framed}
		% \scriptsize
			% \begin{compactdesc}
				% \item[JSP] 
				% % \item[BC] 
				% % \item[]
			% \end{compactdesc}
		% \end{framed}

	% -----
	\subsubsection{Section 66:11} % *DC66-11 
	% -----
		\label{DC66-11}

		Keep these sayings, for they are true and faithful; and thou shalt \hyperref[MAGNIFY]{magnify} thine \hyperref[OFFICE]{office}, and push many people to Zion with songs of everlasting joy upon their heads.

		% \begin{framed}
		% \scriptsize
			% \begin{compactdesc}
				% \item[JSP] 
				% % \item[BC] 
				% % \item[]
			% \end{compactdesc}
		% \end{framed}

	% -----
	\subsubsection{Section 66:12} % *DC66-12 
	% -----
		\label{DC66-12}

		Continue in these things even unto the end, and you shall have a crown of eternal life at the right hand of my Father, who is full%
			\footnote{%
				% \label{}%
				For the phrases ``full of grace and truth'' and just ``grace and truth,'' see \hyperref[TRUTH-PHRASES]{Truth}.
				}
		of grace and truth.%
			\footnote{%
				% \label{}%
				This is the only place in the scriptures which says of the \textit{Father} that he is full of ``grace and truth.'' It's said many times of Christ, or the Son. The only other place that might be saying it of the Father is \hyperref[DC84-102]{DC 84:102}. So in being full of ``grace and truth'' himself, Christ is again just continuing in the attributes the Father has.
				}

		\begin{framed}
		\scriptsize
			\begin{compactdesc}
				\item[JSP] Continue in these things even unto the end and you shall have a Crown of etern[al] life on the right hand of my Father w[ho] is full of grace and truth.
				\item[RB1] continue in those things unto the end \& you shall have a crown of eternal life on the right hand of my father who is full of grace \& truth
				\item[RB2] continue in these things even unto the end and you shall have a crown of eternal life and the right hand of my father who is full of grace and truth
				% \item[BC] 
				% \item[]
			\end{compactdesc}
		\end{framed}

	% -----
	\subsubsection{Section 66:13} % *DC66-13 
	% -----
		\label{DC66-13}

		Verily, thus saith the Lord your God, your Redeemer, even Jesus Christ.  Amen.

		% \begin{framed}
		% \scriptsize
			% \begin{compactdesc}
				% \item[JSP] 
				% % \item[BC] 
				% % \item[]
			% \end{compactdesc}
		% \end{framed}